\documentclass[../main.tex]{subfiles}
\begin{document}
%==========================================TIÊU ĐỀ TẬP========================================
\begin{table}[h]
  \begin{adjustwidth}{-1cm}{}
    \begin{tabular}{>{\centering\arraybackslash}p{6cm}|>{\raggedright\arraybackslash}p{12.5cm}}
      \multirow{5}{6cm}
      % Cột bên trái: ÉPISODE và số tập
      {\centering
        {\\[-40pt]
          \flaregothic\fontsize{20pt}{20pt}\selectfont
          \color{eptype}ÉPISODE\color{black}\\
          \burbank\fontsize{72pt}{72pt}\selectfont
          \color{epnum}86
          \\[8pt]
          \hrule
          \vspace{8pt}
          \flaregothic\fontsize{20pt}{20pt}\selectfont
          \color{eptype}SPÉCIAL\color{black}\\
          \burbank\fontsize{72pt}{72pt}\selectfont
          \color{epnum}XV\color{black}\\[6pt]
          \cafeteria\fontsize{20pt}{20pt}\selectfont\color{epnum}(S-15)\color{black}
          \\[18pt]
          % \flaregothic\fontsize{14pt}{14pt}\selectfont
          % \color{epattr}BIENVENUE, \uline{\color{Banpire} Mel} !
          \flaregothic\fontsize{18pt}{18pt}\selectfont
          \color{holofes}L'HOLOFÊTE ! \\ \color{epattr} {\sen Mini-histoire}\\
          \color{black}
        }
      }
       &
      % Dòng thứ nhất: tên tập tiếng Nhật
      {\fotskip\fontsize{18pt}{18pt}\selectfont これがホロライブの\ruby{年}{ねん}\ruby{末}{まつ}！【\ruby{短}{たん}\ruby{編}{ぺん}\ruby{集}{しゅう}】
      }   \\[6pt]
      % Dòng thứ hai: tên tập tiếng Anh
       & {\cafeteria\fontsize{18pt}{18pt}\selectfont The 2\textsuperscript{nd} Fes: Short Stories}                         \\[4pt]
      % Dòng thứ ba: tên tập tiếng Việt
       & {\vnmsans\fontsize{18pt}{18pt}\selectfont Gặp nhau cuối năm 2020: Chuyện chưa kể trong 2\textsuperscript{nd} Fes} \\[8pt]
      % Dòng thứ tư: Nhân vật xuất hiện
       & {\caviar{\fontsize{14pt}{14pt}\selectfont Nhân vật: Theo từng câu chuyện}}
      \\[8pt]
      % Dòng cuối cùng: Thời gian diễn ra của tập (thường sẽ là ngày phát hành)
       & {\continuum\fontsize{16pt}{16pt}\selectfont Ngày phát hành: 27/12/2020}                                           \\[36pt]
    \end{tabular}
  \end{adjustwidth}
\end{table}
%=========================================TRUYỆN CHÍNH========================================
\color{black}

\vspace{80pt}

\txtbx[2]
{{\color{vol} \uline{Tóm tắt:}} Một tập hợp 10 câu chuyện nhỏ cuối năm đầy hỗn loạn: Marine quảng cáo 2nd Fes bằng sách BL thịt xông khói - rau diếp; Sora bị phỏng vấn với Pekora đang co giật như đồng hồ báo thức; Aqua đổ mù tạt trong phòng xông hơi và Shion làm bánh kẹp từ đó; Miko chụp ảnh từ dưới váy Choco rồi lập trình game hẹn hò; Noel ghen tuông định khắc tượng gỗ Flare; Mio phát điên vì thiếu mèo; Roboco làm sáo thu hút và vô tình hút cả đám đông cùng Matsuri; Suisei rủ Ayame săn khủng long trong khi ba kẻ quậy phá; A-chan và Sora tìm ra bí quyết cho \textit{hologra} nhờ Kuon; và Kuon cùng bạn học xem lại những tập phim xàm xí cả năm.
}

\begin{center}
  \uline{(Tập này được làm nhằm phục vụ cho hậu buổi concert \textit{hololive 2nd fes. Beyond the Stage} vào ngày 21 và 22/12/2020.)}
\end{center}

Trên toàn bộ văn phòng \hll, không khí cuối năm đang lan tỏa khắp mọi ngóc ngách. Ánh đèn từ các phòng làm việc vẫn thắp sáng rực rỡ, những chiếc máy tính chưa bao giờ tắt hẳn để phục vụ cho công việc không ngừng nghỉ của các cô gái thần tượng, nhưng trên môi ai cũng nở nụ cười háo hức chờ đợi sự kiện lớn nhất năm: buổi hòa nhạc \textit{hololive 2nd fes. Beyond the Stage}. 

Những ngày trước sự kiện quan trọng này, tất cả đều bận rộn với lịch trình dày đặc: luyện tập vũ đạo, thu âm các bài hát mới, chuẩn bị trang phục cho buổi biểu diễn. Nhưng giữa sự tất bật đó, những câu chuyện nhỏ vẫn nảy sinh hàng ngày, như những mảnh ghép đầy màu sắc vẽ nên bức tranh về cuộc sống thật sự của các thành viên \hll\ - không chỉ là những thần tượng tỏa sáng trên sân khấu, mà còn là những người bạn đồng hành, cùng trải qua những phút giây điên rồ, đáng yêu và không thể nào quên.

Tập đặc biệt này sẽ tổng hợp 10 câu chuyện ngắn, ghi lại những khoảnh khắc thú vị, hài hước, đôi khi lại kỳ quặc đến khó tin diễn ra trong những ngày cuối năm 2020, trước khi toàn bộ 35 thành viên cùng đứng lên sân khấu lớn để mang đến buổi biểu diễn đáng nhớ trong lịch sử \hll. Mỗi câu chuyện là một mảnh ghép, kể về tình bạn, sự điên rồ và những kỷ niệm không thể nào thay thế của những cô gái đang cùng nhau theo đuổi ước mơ.

\vspace{-8pt}

\subsection*{\phantomsection\label{ep-S13.I}}
\addcontentsline{toc}{subsection}{{\sen{-Histoire 1-}} {\caviar Cuốn sách BL}}
\stpt{-Histoire 1-}\stcap{Cuốn sách BL}
\begin{center}
{\caviar\fontsize{12pt}{12pt}\selectfont Nhân vật: \emp{rushia} \emp{marine}}
\end{center}

Marine chìa ra một cuốn sách với tựa đề đầy nghi vấn, nụ cười nở trên môi như một kẻ vừa tìm thấy báu vật: ``Nè, Rushia\~{} Đọc thử cái này đi.'' Trên bìa cuốn sách, dòng chữ ``Cuốn sách BL - Thịt xông khói x Rau diếp'' được in đậm, toát lên vẻ vừa buồn cười vừa mờ ám.

\img{7-4a1}

Nét mặt nàng Hải tặc như một kẻ dâm đãng chính hiệu, đôi mắt cô long lanh với vẻ khoái trá. Cô sán lại gần Rushia và thì thầm với giọng điệu đầy hứng thú, như thể đang chia sẻ một bí mật thiêng liêng: ``Đây là một cuốn sách quý về chuyện tình lãng mạn giữa thịt xông khói và cây rau diếp đấy\~{}''

Nàng Pháp sư lập tức đổ mồ hôi hột, trán nhăn lại như không thể hiểu nổi những gì vừa nghe. Cô lùi lại một bước, giọng đầy cảnh giác: ``Ờm\dots\ bà đang nói cái quái gì vậy-nanodesu?''

Nhưng Marine chẳng buồn giải thích, cô nhẹ nhàng lướt tay trên bìa sách, đôi mắt mơ màng như thể đang chiêm ngưỡng một tuyệt tác nghệ thuật vô giá. Giọng cô tiếp tục nhỏ nhẹ, nhưng lại khiến người ta rùng mình vì sự nhiệt tình quá đỗi: ``Bà có biết không? Đây chính là sự hòa quyện giữa hai hương vị tuyệt vời. Những lát thịt xông khói dai dai, hòa quyện cùng những lớp rau diếp giòn tan\dots\ Và chúng còn được `tắm' trong lớp nước sốt vừng đầy quyến rũ nữa!''

Rushia nheo mắt đầy hoài nghi, cả người cứng đờ trước những lời lẽ mờ ám của đồng nghiệp. Cô cố gắng vùng vẫy, giọng lắp bắp: ``N-Này! Bỏ tôi ra!''

Nhưng nàng Hải tặc chẳng hề có ý định dừng lại. Cô uốn éo lắc lư cái hông, đôi mắt long lanh như thể đang truyền bá một tôn giáo mới, giọng cô càng lúc càng hưng phấn: ``Bà biết không, Rushia? Những thớ thịt mềm mại của thịt xông khói\dots\ Những lá rau diếp mọng nước, nhai vui tai đến mức kích thích giác quan\dots\ Tất cả hòa quyện trên bàn ăn như một cặp tình nhân nồng cháy!''

\img{7-4a2}

Kuon người đang ngồi cách đó không xa, nghe những lời mật ngọt như vậy, không thể tập trung viết báo cáo được nữa. Tiếng thì thầm đầy ẩn ý của nàng Hải tặc đã xuyên qua mọi rào cản. Cậu từ từ ngẩng mặt lên, nhìn hai thành viên thần tượng đang dây dưa theo nghĩa đen, đầu óc cậu như rối tung. Trong đầu cậu lúc này chỉ có một câu duy nhất: ``\vie{\dots Cái lờ gì thế?}''

Nàng Pháp sư lúc này đã lắp bắp không nói nên lời, mặt đỏ bừng vì ngượng và bối rối: ``Ý-Ý của bà là gì chứ!?''

Bất ngờ, nàng Hải tặc liền đứng thẳng dậy, áp sát thẳng mặt vào cô đồng nghiệp, rồi hét lên một câu như thể vừa nhận được linh cảm từ trên cao: ``\textbf{Buổi \hll\ 2nd Fes - Beyond the Stage! Tài trợ bởi Bushiroad!}''

Cả cậu Tổng và nàng Pháp sư cùng đồng thanh, giọng đầy sự khó hiểu: ``Bà/Em đang nói cái củ kẹc gì thế!?''

Cuối cùng, không ai hiểu vì sao Marine lại hét lên tên chương trình mà họ đang biểu diễn vào giữa một cuộc trò chuyện về thịt xông khói và rau diếp. Có lẽ vì cô sợ khán giả mất tập trung\dots\ hoặc đơn giản là cô đang nghĩ cách quảng bá chương trình theo một phong cách rất riêng của mình.

\newpage

\subsection*{\phantomsection\label{ep-S13.II}}
\addcontentsline{toc}{subsection}{{\sen{-Histoire 2-}} {\caviar Cuộc phỏng vấn kỳ lạ}}
\stpt{-Histoire 2-}\stcap{Cuộc phỏng vấn kỳ lạ}
\begin{center}
{\caviar\fontsize{12pt}{12pt}\selectfont Nhân vật: \emp{kuon1} \emp{sora} \emp{choco} \emp{pekora}}
\end{center}

Tại một góc bàn, Choco và Kuon ngồi đối diện với Sora, sẵn sàng cho buổi phỏng vấn đặc biệt. Ánh đèn vàng từ chiếc đèn bàn hắt xuống, tạo ra một bầu không khí vừa trang trọng vừa kỳ quặc.

Ở giữa họ, Pekora nằm gục xuống bàn như một cái xác không hồn, tay cô buông thõng xuống hai bên. Trên đầu cô nàng thỏ còn áp một cái đồng hồ báo thức, từng cơn co giật nhẹ chạy dọc sống lưng theo từng nhịp như thể tinh thần đã rời khỏi thể xác. Chiếc đồng hồ phát ra tiếng tích tắc đều đặn, hòa cùng nhịp thở yếu ớt của cô.

Nàng Y tá hào hứng tuyên bố, giọng cô vang dội như một phát thanh viên truyền hình: ``Bọn cô sẽ bắt đầu cuộc phỏng vấn nha!''

Nàng Thần tượng nhẹ nhàng gật đầu đồng ý, nụ cười dịu dàng vẫn thường trực trên môi, nhưng chưa kịp chuẩn bị tinh thần thì câu hỏi đầu tiên đã ập đến như một cú sốc. ``Câu hỏi đầu tiên: Em có thích Landolt C\footnote{Công cụ kiểm tra thị lực của một người dưới dạng một tấm bảng có các chữ C, mỗi chữ C sẽ có một vạch từ bốn hướng: lên, xuống, trái, phải. Kích cỡ chữ C sẽ giảm xuống từ phần trên của bảng xuống phía dưới.} không?''

\img{7-4b1}

Cậu Tổng liếc mắt sang nàng Y tá, vẻ mặt đầy nghi hoặc, đôi mày cậu nhíu chặt lại: ``Anh thấy câu hỏi này kỳ quặc thế\dots\ Mà thôi kệ, em thử trả lời anh xem nào?''

Sora hơi giật mình trước câu hỏi khó hiểu, đôi mắt cô mở to với vẻ bối rối, nhưng vẫn ấp úng đáp lại với giọng đầy lưỡng lự: ``C-Chị nghĩ là nó cũng hay ho đấy\dots\ Nhưng mà khó chọn quá\dots''

Choco lập tức tiếp tục với vẻ mặt đầy hứng thú, cô giơ tay chỉ về phía nàng Thỏ đang nằm bất động: ``Vậy giữa Landolt C và Pekora-sama - người đã làm việc tăng ca suốt 1024 ngày, em thích cái nào hơn?''

Pekora giờ đây trông chẳng khác gì một xác sống, toàn thân rũ rượi, con mắt thì trông như mất hết sức sống, chỉ còn tròng trắng. Đồng hồ báo thức trên đầu cô vẫn kêu tí tách đều đặn, còn cô thì co giật nhẹ từng đợt, như thể đã thoát ly khỏi thực tại và lang thang trong một thế giới khác.

\img{7-4b2}

Kuon đổ mồ hôi hột, nhìn cảnh tượng trước mặt mà hỏi nhỏ, giọng cậu như muốn thì thầm để không đánh thức một thực thể nguy hiểm: ``Em ấy trông có ổn không vậy\dots?''

Sora sững sờ nhìn người đồng nghiệp của mình, cảm giác như mình đang bị đặt vào một tình huống sinh tồn khắc nghiệt, nơi mà mọi lựa chọn đều có thể dẫn đến hậu quả không lường trước. Mồ hôi lạnh bắt đầu rịn ra trên trán cô. ``Hả? Ờ thì\dots\ em thích Pekora-chan hơn,'' cô tỏ vẻ thương hại cho nàng Thỏ tội nghiệp, giọng cô như muốn bảo vệ một sinh mệnh đang hấp hối.

Nàng Y tá lập tức nheo mắt đầy ẩn ý, phát ra một tiếng ``Ồ?'' đầy nghi vấn, khiến không khí bỗng chốc căng thẳng hơn. Cậu Tổng lập tức lắc đầu, giọng cậu đầy sự cảnh báo: ``Anh nghĩ đó không phải là lựa chọn phù hợp đâu. Cứ nhìn em ấy đi là biết. Anh cho em chọn lại đấy.''

Nàng Thần tượng câm nín nhìn nàng Thỏ co giật liên tục trên bàn với những tiếng tích tắc phát ra từ chiếc đồng hồ. Sau một hồi suy nghĩ, và một chút sợ hãi hiện rõ trên gương mặt, cô đành gãi đầu, cười trừ rồi ``quay xe'' một cách dứt khoát, giọng cô như một lời thú nhận đầy áy náy: ``Thôi thì\dots\ chị chọn Landolt C vậy!''

Ngay lập tức, cả nàng Y tá và cậu Tổng đồng thanh hô vang với giọng đầy phấn khích, như thể vừa kết thúc một trò chơi truyền hình: ``Chị/Em được nhận!''

Bất ngờ thay, ngay khoảnh khắc đó, nàng Thỏ như bị đánh thức bởi một thế lực vô hình, cô bật dậy như lò xo, mắt mở to, rồi hét to với giọng đầy máy móc, như thể cô đã trở thành chính chiếc đồng hồ trên đầu: ``5:50 sáng rồi! 5:50 sáng rồi!'' Động tác của cô không khác gì một chiếc đồng hồ báo thức đang vang lên inh ỏi, đầu cô lắc qua lắc lại như còi báo động.

\img{7-4b3}

Vừa hét, cô vừa tự đập vào đầu mình, đôi tai cô cuộn tròn lại làm hai tai đồng hồ, trông như thể đang bắt chước một cái đồng hồ thật, tạo nên một cảnh tượng vừa kỳ lạ vừa đáng thương. Sora tròn mắt kinh ngạc, giọng cô đầy vẻ không tin nổi: ``Chị không biết là Pekora-chan lại là người thích dậy sớm đấy!''

Kuon thở dài, khoanh tay cười chát với vẻ mặt của một người đã chứng kiến quá nhiều điều kỳ quái: ``Rồi em sẽ phải quen với môi trường ở đây thôi. Chào mừng em đến với công ty.''

Còn nàng Y tá thì vỗ tay không ngớt, mặt đầy vẻ thích thú như thể vừa chứng kiến một màn kịch đầy bất ngờ, đôi mắt cô sáng lên với vẻ thích thú: ``Xem chừng có vẻ thú vị đây\dots''

\newpage

\subsection*{\phantomsection\label{ep-S13.III}}
\addcontentsline{toc}{subsection}{{\sen{-Histoire 3-}} {\caviar Xông hơi và\dots\ mù tạt?}}
\stpt{-Histoire 3-}\stcap{Xông hơi và\dots\ mù tạt?}
\begin{center}
{\caviar\fontsize{12pt}{12pt}\selectfont Nhân vật: \emp{kuon1} \emp{aqua} \emp{shion}}
\end{center}

Tại một góc hành lang, Aqua đứng chờ sẵn, vừa thấy bóng cô bạn Shion xuất hiện, cô nàng đã vẫy tay gọi với giọng đầy hào hứng: ``Shion-chan, ta đi xông hơi đi!''

Bước cạnh nàng Phù thủy là Kuon, người đang rảnh rang không có việc gì làm, cậu bước theo với vẻ mặt hờ hững. Nàng Phù thủy nghiêng đầu nhìn cậu, giọng có chút ngạc nhiên: ``Với Kuon-senpai sao?''

Cậu Tổng khoanh tay, nhíu mày đầy nghi hoặc, đôi mắt cậu híp lại với vẻ bất mãn: ``Sao anh có cảm giác không được chào đón ở đây nhỉ\dots?''

Shion lập tức xua tay với vẻ mặt hốt hoảng, như thể vừa bị bắt quả tang: ``E-Em có nghĩ thế đâu!'' rồi cô nàng quay sang cô bạn thân với vẻ mặt đầy quyết tâm: ``Mà, đi xông hơi à? Đi luôn!''

Kuon nhìn cô với ánh mắt khó hiểu, giọng cậu như muốn hỏi về sự thay đổi thái độ đột ngột: ``Chấp nhận dễ dàng đến vậy sao\dots?''

\begin{center}
\uline{Bên trong phòng xông hơi.}
\end{center}

Cả ba ngồi yên trong bầu không khí nóng bức, những làn hơi nước bốc lên nghi ngút bao phủ xung quanh, mồ hôi lấm tấm trên trán họ. Họ vẫn còn đang mặc quần áo trên người, không hề muốn cởi bỏ ra vì sợ đối phương nhìn thấy, tạo nên một khung cảnh vừa buồn cười vừa ngượng ngùng. Cặp đôi ``Hành - Tỏi'' ngồi ở một phía, còn cậu Tổng ngồi ở phía đối diện, tạo thành một tam giác khó xử.

\img{7-4c1}

Shion và Kuon đã bắt đầu thở dốc, mồ hôi chảy như thác nước trên gương mặt họ, những giọt mồ hôi lăn dài xuống cằm. Kuon lau trán, than thở với giọng đầy mệt mỏi: ``Nóng thật\dots''

Nàng Phù thủy cũng không khá hơn, cô cố cử động nhưng cảm giác như toàn thân mình đang tan chảy, giọng cô yếu ớt như sắp đuối sức: ``Em đang rã mồ hôi như thác nước này\dots''

Cậu Tổng nhìn sang Aqua, thấy cô nàng vẫn ngồi yên bình, không hề có dấu hiệu mệt mỏi, mặt cô vẫn tỉnh táo như thể đang ngồi trong phòng điều hòa.

Cậu bĩu môi than thở, tay đưa lên xoa cổ: ``Ai bảo mặc quần áo vào làm gì\dots'' để rồi cậu sực nhớ ra cậu đang không ngồi một mình, cậu quay sang nàng Hầu gái: ``Mà, đây là phòng ghép mà nhỉ?''

Nàng Hầu gái lười biếng trả lời, giọng cô vẫn bình thản như không có gì xảy ra: ``Hết phòng riêng rồi nên mới vậy chứ\dots'' Nàng Phù thủy cười nhẹ, thả lỏng người hơn, cảm nhận hơi nóng đang xoa dịu cơ thể: ``Xông hơi đúng là phê thật, đúng không nhỉ, Aqua\dots''

Cô vừa quay sang bên cạnh, định trò chuyện với cô bạn thân của mình thì bất chợt nhìn thấy có hiện tượng lạ trên người nàng cô bạn. Trên cơ thể nàng Cô hầu, những thứ chất lỏng đang chảy xuống từng giọt, nhưng thay vì là mồ hôi trong suốt, nó có một màu vàng kỳ quái, như thể có thứ gì đó khác đang thoát ra từ cơ thể cô.

\img{7-4c2}

Shion trợn tròn mắt, giọng cô đầy sự kinh ngạc: ``Cậu ta đang đổ mồ hôi trộm\dots\ à không, là thác Niagara mới đúng!''

Kuon nhìn một lúc rồi hốt hoảng, cậu nhận ra màu sắc quen thuộc đó, một màu vàng đặc trưng mà bất kỳ ai cũng biết: ``\dots Không, anh thấy em ấy đang đổ mù tạt ra đấy!''

Cả hai người lập tức lo lắng, cảm giác không khí trong phòng càng thêm ngột ngạt, như thể họ đang ngồi trong một nồi lẩu khổng lồ.

Nàng Phù thủy thở dài, lắc đầu bất lực với vẻ mặt của một người đã quá quen với những điều kỳ quặc của bạn mình: ``Trời ạ, Aqua-chan, tôi đã bảo bao nhiêu lần là đừng có giấu mù tạt trong người nữa mà!''

Kuon nhìn cô nàng hầu gái với ánh mắt khó hiểu, đầu cậu nghiêng sang một bên như một chú chó đang cố gắng hiểu mệnh lệnh: ``Để thế làm chua người chứ làm cái gì hả, Hành tím?''

Trong lúc cậu Tổng còn đang tìm cách xử lý chỗ mù tạt mà nàng Hầu gái tiết ra, nàng Phù thủy chợt nảy ra một ý tưởng với nụ cười tinh quái: ``A, nói đến hành tím thì\dots'' Cô nhếch môi cười tinh quái, mắt sáng rực như vừa tìm ra công thức vĩ đại: ``Nếu ta kẹp cậu ta với bánh mì kẹp thịt xông khói và rau diếp thì sao nhỉ?''

Cô nàng lấy ra một chiếc bánh mì kẹp cất sẵn trong người, tỏa ra mùi thơm của thịt nướng và rau tươi. Kuon lập tức quay ngoắt sang nhìn cô với đôi mắt mở to, giọng đầy sửng sốt: ``Em kiếm cái đó ở đâu ra vậy?''

\img{7-4c3}

Rồi cậu nhăn mày, một cảm giác quen thuộc dâng lên trong đầu, như thể vừa nghe thấy một điều gì đó đã từng xuất hiện trong một câu chuyện khác: ``Khoan đã\dots\ thịt xông khói với rau diếp á? Sao nghe quen vậy?''

Shion không trả lời, chỉ lẳng lặng lấy chiếc bánh kẹp ra, đưa nó vào dòng mồ hôi mù tạt đang chảy từ nàng Hầu gái, hứng đầy lớp sốt màu vàng lên bề mặt bánh rồi cắn một miếng to! Cô nàng lập tức reo lên với vẻ mặt đầy thỏa mãn: ``Ngon quá!''

Kuon sững người nhìn cô đang ăn miếng bánh ngon lành, giọng cậu đầy vẻ không tin nổi: ``\vie{Vãi.} Vậy cũng được nữa\dots?''

Nhưng rồi thấy cô nàng đang hướng chiếc bánh kẹp về phía cậu, cậu cũng tỏ vẻ tò mò, mùi thơm của bánh mì và sốt mù tạt tươi mới đang kích thích vị giác: ``Này, anh thử một miếng được không?''

\newpage

\subsection*{\phantomsection\label{ep-S13.IV}}
\addcontentsline{toc}{subsection}{{\sen{-Histoire 4-}} {\caviar Nghệ thuật vì đam mê}}
\stpt{-Histoire 4-}\stcap{Nghệ thuật vì đam mê}
\begin{center}
{\caviar\fontsize{12pt}{12pt}\selectfont Nhân vật: \emp{kuon1} \emp{ryuuhou1} \emp{miko2} \emp{mel} \emp{choco} \emp{pekora}}
\end{center}

Một buổi chiều mát mẻ lộng gió, Miko đang đi dạo phố - một điều mà một người hay giam mình trong phòng rất hiếm khi làm. Thay vì ngồi trong phòng với những thiết bị điện tử, cô bước ra ngoài với vẻ mặt đầy bí ẩn. Theo sau cô là Mel, Pekora, Kuon và Ryuuhou, người đang vừa đi vừa lướt điện thoại.

\img{7-4d1}

Nàng Ma cà rồng nghiêng đầu thắc mắc, đôi mắt cô lấp lánh với vẻ tò mò: ``Không biết tại sao tự dưng Miko-chan lại muốn ra ngoài nhỉ?''

Nàng Thỏ chống cằm suy nghĩ, đôi tai vểnh lên như hai cái ăng-ten dò tìm thông tin: ``Pekora cũng tò mò lắm ha, peko.''

Cậu Tổng nhún vai, vẻ mặt thản nhiên với một nụ cười nhẹ: ``Anh lại thấy bình thường mà. Giam mình mãi trong phòng cũng tù túng, thỉnh thoảng cũng nên ra ngoài hít thở tí chứ.''

Cô em gái của cậu ngước lên khỏi điện thoại, thêm vào: ``Em nghĩ chị ấy có mục đích gì đó bí mật lắm, kiểu như đi săn ảnh đẹp ấy mà.''

Bất chợt một cơn gió mạnh thổi qua, làm mọi người giật mình, tóc và váy của các cô gái tung bay trong không trung. Nhưng người ngạc nhiên nhất lại là Choco - cách họ phía trước vài bước chân - vì cơn gió vô tình hất tung váy của cô lên, để lộ một khoảnh khắc mà cô không mong muốn.

``Kya\~{}!'' nàng Y tá thét lên, vội vàng giữ váy lại với khuôn mặt đỏ bừng.

Nàng Thỏ ôm đầu, đôi tai rũ xuống vì bất ngờ: ``Gió mạnh quá, peko!''

Nàng Trưởng nhóm với bộ trang phục không có váy, chỉ đứng đó cười khúc khích, vân vê mái tóc. Cô liếc nhìn anh trai mình với ánh mắt tinh nghịch: ``Ối giời, may mà em không mặc váy hôm nay. Không thì anh lại được ``hưởng lợi'' rồi còn gì.''

Cậu Tổng nhướng mày, đáp lại với giọng đầy vẻ khinh khỉnh: ``Ai thèm chứ. Em nghĩ anh rảnh đến thế à?'' Ryuuhou cười hì hì, đáp lại: ``Nói thế chứ mắt anh nhìn chỗ nào suốt lúc gió thổi vậy?''

Kuon lắc đầu, không thèm đáp lại em gái, mà ngước lên bầu trời, khẽ nhíu mày khi nhìn những đám mây đang kéo đến: ``Kiểu này trời sắp mưa rồi đấy\dots\ Lại không mang ô mới đau chứ.''

Nàng Ma cà rồng đưa tay vuốt tóc, lắc đầu cảm thán với giọng đầy triết lý: ``Những cơn gió mạnh bất chợt thế này quả là bất tiện cho những người mặc váy nhỉ.''

Ryuuhou chen vào với vẻ mặt hả hê: ``Thế nên em mới chọn áo phông với quần thể dục đấy. Tiện lợi, thoải mái, lại không lo bị lộ.''

Kuon gật gù với vẻ mặt như một nhà quan sát đã chứng kiến quá nhiều: ``Ừ. Cơ mà nó lại là cơ hội cho những kẻ biến thái\dots\ như em ấy chẳng hạn.'' Cậu hất cằm về phía Miko, người hiện đang đeo kính râm, trên tay cầm một chiếc máy ảnh, lia ống kính khắp nơi với vẻ mặt tập trung cao độ. Nhưng thay vì chụp cảnh vật, nàng Trông đền đang tập trung chụp vào những góc\dots\ rất khó diễn tả.

\gif{7-4d2}{1}{8}

Ryuuhou tròn mắt khi thấy những gì Miko đang làm, cô bé thì thầm với anh trai: ``Onii-san, Miko-san đang chụp ảnh từ dưới váy Choco-sensei kìa. Trời ơi, không ngờ chị ấy lại dữ vậy.''

Pekora vội kéo tay Mel, giọng cô như một phóng viên đưa tin nóng: ``Hai người thấy chưa!? Miko-chi trong nháy mắt đã lôi cái máy ảnh ra và chụp rất nhiều pô từ hướng lên trên của váy Choco-sensei kìa, peko!''

Nàng Dơi há hốc miệng, giọng đầy sửng sốt: ``Ể!? Thật á!?'' Kuon khoanh tay, gật gù đầy cảm khái, giọng như một nhà phê bình nghệ thuật đang bình luận về một tác phẩm vĩ đại: ``Nhìn cái cách lăn xả của em ấy là biết em ấy hy sinh vì nghệ thuật thế nào rồi.''

Ryuuhou lắc đầu với vẻ vừa buồn cười vừa bất lực: ``Nghệ thuật gì mà nhìn khiêu dâm quá vậy anh.''

Mel trợn mắt, giận dữ lay vai cậu, giọng cô đầy sự bất bình: ``Nghệ thuật cái gì chứ! Anh không định ngăn Miko-chan kia sao!?''

Kuon thở dài, giọng đầy bất lực của một người đã từ bỏ hy vọng: ``Mấy đứa biết não Miko-chi chỉ toàn có eroge với mấy cái khiêu dâm thôi. Ngăn thế đếch nào được.''

''Cái gì chứ!?'' nàng Ma cà rồng và nàng Thỏ thảng thốt, không thể tin nổi với những gì họ nghe, mặt họ tái nhợt vì sốc.

Ryuuhou cũng trợn mắt nhìn anh trai: ``Anh nói thật á? Miko-san mà làm eroge cơ đấy?''

``Cũng không hẳn,'' ông anh trai đáp, mắt liếc về phía nàng Vu nữ vẫn đang lăn xả cầm máy ảnh, ``em ấy đã dâm đãng như thế từ khi còn chưa vào \hll\ kia mà. Thậm chí còn từng stream eroge trên kênh của mình nữa cơ, may mà toàn là bản mọi lứa tuổi.''

Pekora vỗ trán, giọng vừa ngán ngẩm vừa thảng thốt khi sự bất ngờ từ nàng Vu nữ chưa dừng lại ở việc chụp ảnh: ``Giờ chị ấy còn đang lập trình một tựa game dating sim của chính tay chị làm kìa, peko!!!''

\img{7-4d3}

Kuon liếc sang nàng Vu nữ, nhướng mày hỏi với giọng đầy hoài nghi của một lập trình viên: ``Nhưng mà có biết viết code không đấy? Trình độ ê-lít như em ấy còn chưa dùng nổi VS nữa\dots''

Nàng Vu nữ không trả lời, cô chỉ cười như một tên gàn dở, nụ cười nở rộ trên gương mặt với đôi mắt đầy mơ màng. Trong đầu cô giờ đây chỉ có những cảnh nhạy cảm, và bản thân cô lại đang chìm đắm vào sở thích ``quái thai dị dạng'' này, không một chút hối hận.

Cậu Tổng nhấp một ngụm nước, bình thản hỏi với giọng điềm tĩnh của một người đã quá quen: ``Có muốn anh phụ phần lập trình không?''

Ngay lập tức, Miko gật đầu lia lịa, đôi mắt cô sáng rực lên như hai ngọn đèn pha. Không hiểu bằng cách nào một chiếc ghế đột nhiên xuất hiện ngay bên cạnh cô. Cậu thản nhiên ngồi xuống, bắt đầu thảo luận với nàng Vu nữ về trò chơi.

Ryuuhou nhìn cảnh tượng đó, cô bé khoanh tay và nói với giọng đầy nghi ngờ: ``Onii-san, anh nói là chỉ phụ giúp code thôi đúng không? Chứ không phải phần hình ảnh nhé?''

Kuon đáp lại với giọng bình thản: ``Ừ, chỉ code thôi. Anh không dính vào phần hình ảnh gì hết.''

Cô em gái nheo mắt, giọng đầy vẻ không tin: ``Anh nói thế em không tin đâu. Ai mà biết được anh với Miko-san đang bàn bạc cái gì. Em thấy hai người có vẻ khá là hào hứng đấy.''

``Em đừng có suy diễn lung tung. Anh chỉ giúp đỡ thôi.''

''Này, này! Cho Pekora tham gia với!'' Pekora cũng không nhịn được, sà vào hai người họ đang bắt đầu nói chuyện rất rôm rả, đôi tai thỏ dựng lên vì hứng thú.

Ryuuhou nhìn cảnh tượng ba người đang cặm cụi bàn bạc về code và game, cô bé lắc đầu với vẻ mặt đầy bất lực, thì thầm với nàng Ma cà rồng: ``Mel-san à, em thấy mấy người đó không còn cứu vãn được nữa rồi.''

Nàng Ma cà rồng nhìn cảnh tượng trước mắt, biểu cảm pha trộn giữa ngỡ ngàng, giận dữ, chán chường và một chút kinh tởm. Cô chỉ biết đứng đó, lắc đầu bất lực trước những gì đang diễn ra xung quanh mình.

\newpage

\subsection*{\phantomsection\label{ep-S13.V}}
\addcontentsline{toc}{subsection}{{\sen{-Histoire 5-}} {\caviar Nghệ thuật điêu khắc vì tình yêu}}
\stpt{-Histoire 5-}\stcap{Nghệ thuật điêu khắc vì tình yêu}
\begin{center}
{\caviar\fontsize{12pt}{12pt}\selectfont Nhân vật: \emp{kuon1} \emp{flare1} \emp{rushia} \emp{noel}}
\end{center}

Trong khu rừng vắng, Flare và Rushia đang đi dạo, vừa đi vừa trò chuyện với những tiếng cười vui vẻ vang vọng giữa tán cây xanh mướt. Nàng Yêu tinh vàng cười đùa với tai nạn do mình gặp phải vừa rồi, giọng cô đầy vẻ hài hước: ``Nó dính đến mức mà tôi không thể thở được à!''

Nàng Pháp sư gật gù, tò mò hỏi với vẻ mặt đầy thích thú: ``Thế à\dots\ vậy bà đã thoát ra bằng cách nào vậy?''

Đằng sau họ, một bóng người đang ẩn nấp trong bụi cây - Noel, với khuôn mặt đầy ganh tị nhìn cảnh hai người bạn thân đang vui vẻ trò chuyện. Kuon đứng phía sau thì khoanh tay, cười khành khạch với vẻ mặt của một khán giả đang thưởng thức một vở kịch hay.

\img{7-4e1}

Nàng Hiệp sĩ nghiến răng, mắt long sòng sọc, tay run run nắm lấy thân cây gần đó với vẻ mặt đầy ghen tức: ``Flare\dots\ Tại sao bà có thể vui vẻ nói chuyện với người khác ngoài tôi chứ?''

Cậu Tổng ở phía sau liền hắng giọng, rồi đột nhiên cất tiếng hát từ bài của cặp đôi đã hát phiên bản thứ hai của một bài hát nhân dịp ``Rô-cô-na'', giọng cậu đầy nhại: ``Bởi vì em ghen, ghen, ghen, ghen mà\~{}''

Câu hát của cậu như đổ thêm dầu vào lửa, khiến cô nàng càng tức tối hơn, khuôn mặt cô đỏ bừng vì ghen tuông. Cô lôi ra một cây búa, mặt hằm hằm với đôi mắt sáng quắc: ``Giờ thì tôi tức rồi đấy!''

Kuon giật mình, lùi lại một bước với vẻ mặt đầy cảnh giác, tay cậu giơ lên như một bức tường phòng thủ: ``Thôi, bỏ búa xuống đi. Kẻo lại làm hại người ta bây giờ!''

Noel lườm cậu, thở dài một hơi đầy bực bội, hạ cánh tay đang giơ búa xuống: ``Em có định đánh cậu ta đâu.'' Kuon ngơ ngác, giọng cậu đầy vẻ khó hiểu: ``Hả?''

Cô hếch cằm đầy kiêu hãnh, rồi giơ lên một cái đục sáng loáng với ánh mắt đầy quyết tâm: ``Em định chỉ khắc Flare làm từ gỗ cho bản thân em thôi.''

Cậu Tổng đứng hình vài giây, sau đó phì cười với vẻ mặt vừa bất ngờ vừa thích thú: ``Tất nhiên là không thể rồi. Tính cách của em ấy hiền như hoa thế mà\dots\ Cùng lắm cũng chỉ là vô tình gây hại thôi.''

Nàng Hiệp sĩ cúi xuống, lẩm bẩm với ánh mắt sáng rực như một kẻ đang theo đuổi lý tưởng: ``Heh\dots\ Rồi chúng ta sẽ cùng đi tận hưởng tuần trăng mật với nhau sau khi tôi khắc bà xong nha, Flare\~{}''

Nụ cười mất nhân tính của cô khiến cậu Tổng chỉ biết cười xòa, tự nhủ tốt nhất không nên xen vào nữa. Cậu lùi lại thêm vài bước để quan sát từ khoảng cách an toàn. Nhưng ngay khi Noel định vung búa chặt cái cây đầu tiên để bắt đầu tác phẩm điêu khắc, một giọng nói quen thuộc vang lên từ phía xa, như một tiếng chuông cảnh tỉnh: ``Kia chả phải Noel à! Bà đang làm gì đấy?''

Nàng Hiệp sĩ giật bắn người, búa suýt rơi khỏi tay, mặt cô tái nhợt đi vì bối rối: ``F-Flare!?''

\img{7-4e2}

Kuon khoanh tay, nửa đùa nửa thật với một nụ cười hài hước: ``Cầu được ước thấy nha.''

Thấy thế, nàng Hiệp sĩ lúng túng, ấp úng, mặt đỏ bừng vì không ngờ cô bạn thân của mình sẽ tới. Cô hoa tay múa chân, cố tìm một lời bao biện trong hành động dường như vô hại của mình, giọng cô như một kẻ bị bắt quả tang đang cố chạy trốn: ``K-Không phải như bà nghĩ đâu! Chỉ là\dots\ tôi muốn khắc họa thân bà thôi!''

\gif{7-4e3}{1}{13}

Cậu Tổng quản gật gù, nhún vai tiếp lời với giọng đầy vẻ hài hước: ``Ít nhất là những gì em ấy muốn làm, Flare-chan.''

Nàng Yêu tinh vàng nghiêng đầu khó hiểu, đôi mắt cô mở to với vẻ ngạc nhiên: ``Ủa? Em sao? Ơ, thế là thế nào? Anh nói rõ hơn được không?''

Không khí bỗng chốc trở nên ngượng ngùng khi Noel cúi gằm mặt, còn cậu Tổng thì bật cười trước cảnh tượng, để lại Flare đứng đó với đầy câu hỏi không lời đáp.

\newpage

\subsection*{\phantomsection\label{ep-S13.VI}}
\addcontentsline{toc}{subsection}{{\sen{-Histoire 6-}} {\caviar Liều thuốc mèo}}
\stpt{-Histoire 6-}\stcap{Liều thuốc mèo}
\begin{center}
{\caviar\fontsize{12pt}{12pt}\selectfont Nhân vật: \emp{kuon1} \emp{okayu} \emp{korone} \emp{mio}}
\end{center}

Tại văn phòng, Mio bất chợt ôm đầu, mặt xanh lè như tàu lá chuối, đôi mắt chỉ còn có những hình xoáy tròn như thể đang phát điên, cô gào lên với giọng đầy tuyệt vọng: ``\textbf{Á! Mình không có đủ mèo rồi! Mình cần một liều điều trị về mèo thôi!}''

Kuon đứng bên cạnh, mồ hôi túa ra trên trán, giọng cậu đầy bối rối: ``Anh chả hiểu ý em nói là gì, nhưng văn phòng mình làm gì có mèo đâu? Hay là bản năng Sói nguyên thủy của em lại lên cơn rồi?''

Vừa dứt lời, một giọng nói quen thuộc vang lên với vẻ đầy tự tin: ``Em có một bài dã thú đây!''

\img{7-4f1}

Okayu xuất hiện từ cửa, khuôn mặt vui thú nhưng ánh mắt lóe lên sự tinh quái, đôi tai cô dựng lên với vẻ hào hứng. Cậu Tổng khẽ giật mình, rồi vỗ trán với một nụ cười bất lực: ``À, anh quên mất là ta có một đứa Mèo ở đây.''

Cậu quay sang nàng Miêu, tò mò hỏi với giọng đầy nghi ngờ: ``Thế em định làm gì để `chữa bệnh' cho Sói đen-chan đây?''

Cô nàng không nói gì, chỉ lẳng lặng đẩy vào trong một nàng Khuyển đang xoay vòng vòng như một chiếc ghế xoay, với cô nàng đang nằm ngửa, hai chân đá lên cao và hai tay co quặp lại, tạo thành một vòng xoáy lông xù.

\gif{7-4f2}{1}{27}

``\textbf{WowWowWowWowWowWowWowWow!}''

Kuon nhíu mày, chớp mắt vài lần trước cảnh tượng quái dị trước mặt, tay cậu đưa lên cằm với vẻ suy tư: ``Thế quái nào Koro-chan có thể quay được mà không thấy bị chóng mặt nhỉ? Mà đúng hơn là em ấy xoay từ khi nào và bằng cách nào vậy?''

Nàng Mèo dường như bơ luôn câu hỏi của cậu, mà chỉ đứng khoanh tay, nghiêm túc nói với giọng của một bác sĩ đang chỉ định phương pháp điều trị: ``Đừng có trốn tránh nữa, Koro-san!''

Nàng Khuyển tiếp tục xoay vòng vòng, hét lên với giọng vừa phấn khích vừa bất lực: ``\textbf{Tôi không thể dừng lại được!}''

Cô cứ thế quay thêm mấy vòng nữa, nhưng rồi chậm dần, chậm dần\dots\ và cuối cùng dừng hẳn, không hề tỏ vẻ lảo đảo gì. Nàng Mèo tím thấy thế thì thở phào nhẹ nhõm lau mồ hôi, giọng đầy hài lòng: ``Thế là xong.''

Nhưng cậu Tổng thì vẫn lắc đầu với vẻ mặt của một người đã chứng kiến quá nhiều điều kỳ lạ: ``Ờm, anh nghĩ em chưa xong đâu.''

Nàng Mèo nghiêng đầu, chưa kịp phản ứng thì bất ngờ nàng Sói đen bỗng nhảy vào, hét lên với giọng đầy phấn khích: ``\textbf{Đây là cách điều trị mèo đó!!}'' rồi bắt đầu điên cuồng nhảy múa như thể vừa trúng bùa độc gì đó. Trông cô nàng chẳng khác gì một người vừa hút xì ke ma túy, hay một người đang bị phát dại, tay chân vung vẩy loạn xạ.

\img{7-4f3}

Cả Kuon và Okayu đều im lặng nhìn cảnh tượng hỗn loạn trước mặt, miệng há hốc, không biết phải phản ứng thế nào.

Cuối cùng cậu Tổng chán nản thở dài, đưa tay xoa trán với vẻ bất lực: ``Anh nghĩ em ấy bị phê đá rồi\dots''

\newpage

\subsection*{\phantomsection\label{ep-S13.VII}}
\addcontentsline{toc}{subsection}{{\sen{-Histoire 7-}} {\caviar Cánh tay máy thu hút}}
\stpt{-Histoire 7-}\stcap{Cánh tay máy thu hút}
\begin{center}
{\caviar\fontsize{12pt}{12pt}\selectfont Nhân vật: \emp{kuon1} \emp{roboco2} \emp{aki} \emp{matsuri}}
\end{center}

Roboco đang hì hụi dưới sàn, tự mình điều chỉnh lại một bộ phận trong cánh tay với vẻ tập trung của một kỹ sư đang hoàn thiện cỗ máy cuối cùng. Ở phía sau, Aki đứng khoanh tay quan sát với ánh mắt tò mò, còn Kuon thì nhìn với vẻ thận trọng, đôi mắt cậu híp lại như đang đánh giá một thiết bị nguy hiểm.

``Roboco-san, chị sửa xong cánh tay chưa thế?'' nàng Dị giới cúi xuống, nhìn cánh tay đang được điều chỉnh mà tò mò hỏi, giọng cô đầy háo hức.

``Rồi, chị đã chỉnh nó thành một cây sáo thu hút rồi,'' nàng Người máy đáp lại, hoàn thành những công đoạn cuối cùng với vẻ mặt tự hào, tay cô vuốt nhẹ cánh tay vừa được nâng cấp.

\img{7-4g1}

``Lần này em lại bày trò mèo gì nữa đây?'' cậu Tổng nhướng mày tỏ vẻ nghi hoặc, tay cậu chống hông. ``Anh hy vọng là em không bày trò rắc rối gì nữa\dots''

``Công nghệ của chị lúc nào cũng thú vị thật đấy! Nhưng nó hoạt động kiểu gì vậy?'' nàng Tóc bay hỏi, đôi mắt sáng lên với vẻ thích thú của một nhà thám hiểm, bỏ ngoài tai hoàn toàn lời cảnh báo của Kuon.

Nàng Người máy cười hì hì rồi giơ cánh tay lên, bên trong cánh tay giờ có một cơ chế nhìn như một cây sáo tinh xảo, với những lỗ nhỏ và đường nét trau chuốt: ``Chị có thể phát ra sóng thu hút từ đây. Chỉ cần thổi một tiếng, đối tượng sẽ bị mê hoặc ngay lập tức!''

Cậu Tổng nhướng mày, nhìn cô với ánh mắt không mấy tin tưởng, giọng cậu đầy hoài nghi: ``Cái này có thực sự cần thiết không?''

``Tất nhiên rồi!'' cô hí hửng đáp, đôi mắt cô sáng rực. ``Thử nghĩ xem, nếu chị rơi vào tình huống nguy hiểm mà không thể di chuyển, chỉ cần thổi sáo, ai đó sẽ chạy đến giúp ngay lập tức!''

``Nghe cũng hợp lý đấy!'' Aki đáp lại một cách ngây thơ, choáng ngợp với cây sáo này, tay cô đưa lên miệng như đang hình dung ra cảnh tượng.

``Nghe hợp lý ở đâu vậy, Aki-chan\dots'' cậu Tổng lắc đầu, tặc lưỡi với vẻ của một người đã có quá nhiều kinh nghiệm với những phát minh của nàng Rô-bốt. ``Thế em không tính đến chuyện nhỡ đâu kẻ được chạy đến lại là một tay biến thái thì sao?''

Nàng Người máy hí hửng chuẩn bị thử nghiệm, nhưng vì bản tính ngốc nghếch sẵn có, cô vô tình kích hoạt chế độ thu hút mà không hề hay biết. Một luồng sóng vô hình phát ra từ cánh tay cô, lan tỏa khắp hành lang.

Vài giây sau, từ hành lang, hàng loạt người không mặt không mũi, toàn thân chỉ có một màu đen xì bất thình lình đổ xô vào căn phòng, tạo thành một dòng người đông đúc như thủy triều. Cả nàng Tóc bay và cậu Tổng đều ngơ ngác khi chứng kiến cảnh tượng trước mặt, miệng há hốc vì kinh ngạc.

\img{7-4g2}

``Á! Một nhóm học sinh không thể kiềm chế được ham muốn bản thân đang lao về phía cánh tay chị ấy kìa!\footnote{Nguyên văn: ​リビドーを抑えきれない学生が集まってきたよ！ (tạm dịch: Có nhiều học sinh không thể kiềm chế được ham muốn tình dục đang lao vào tìm cánh tay của Roboco kìa!)}'' nàng Dị giới thảng thốt, tay cô đưa lên miệng với vẻ kinh hãi. Càng lúc lượng người như vậy càng đông, khiến văn phòng dần dần chật ních như một buổi hòa nhạc đông nghẹt.

``\dots Và lại PON như mọi khi nhỉ, Đại PON-chan?'' cậu Tổng cười khẩy, ngán ngẩm nhìn dòng người đen kịt đang bâu lấy cô nàng Người máy đang ngơ ngác kia, tay cậu đưa lên che mặt vì bất lực.

Roboco hoảng hốt khi thấy mình bị vây quanh bởi một đám người có những hành động quái dị. Họ bắt đầu nhảy xung quanh cô một cách kỳ lạ, tạo thành một vòng tròn như đang thực hiện nghi thức gì đó, đôi tay họ vươn về phía cánh tay của cô với vẻ mê muội.

``\textbf{Em không có ngốc mà!!}'' nàng Rô-bốt đáp lại trong những tiếng hò reo của đám người kia, giọng cô vừa phẫn nộ vừa tuyệt vọng.

Trong lúc cả ba còn chưa kịp tìm cách xử lý, từ phía đám đông, một gương mặt quen thuộc bất ngờ xuất hiện. Một cô gái với nụ cười gian xảo đang chầm chậm tiến lại gần với đôi mắt sáng rực đầy hưng phấn, như một kẻ săn mồi đã tìm thấy con mồi.

\img{7-4g3}

``Matsuri-chan!?'' cả ba người chết sững trong giây lát, không ai tin vào mắt mình.

``Không thể nào\dots\ Cái này mà cũng tác dụng đến cả `quân' của mình ư!?'' Kuon thảng thốt, giọng cậu vỡ òa vì không thể tin nổi, mặt cậu tái nhợt đi.

Aki bàng hoàng, tay cô run lên: ``Ai đó nhanh chóng giải tán nhóm người này đi với!''

``A-Aki-chan, Kuon-senpai, giúp chị với!!!'' Roboco vọng lại, vừa vật lộn với đám người nhưng không ra người ngợm, vừa la hét trong tuyệt vọng, tay cô đưa lên che đầu.

Matsuri tiến thêm một bước, nụ cười càng lúc càng ranh mãnh hơn, đôi mắt cô như hai ngọn đèn pha đang chiếu thẳng vào mục tiêu. Đám đông sau lưng cô cũng đang chờ đợi, mà chắc chỉ có trời mới biết chúng sắp làm gì tiếp theo.

\newpage

\subsection*{\phantomsection\label{ep-S13.VIII}}
\addcontentsline{toc}{subsection}{{\sen{-Histoire 8-}} {\caviar Tên gọi}}
\stpt{-Histoire 8-}\stcap{Tên gọi}
\begin{center}
{\caviar\fontsize{12pt}{12pt}\selectfont Nhân vật: \emp{kuon1} \emp{suisei} \emp{fubuki} \emp{matsuri} \emp{haato} \emp{subaru} \emp{ayame}}
\end{center}

Suisei gọi Ayame, người vẫn đang chăm chú vào bài tập, không buồn ngẩng lên dù chỉ một chút. Bên cạnh cô, Kuon đang ngồi khoanh tay, mắt lơ đãng nhìn về phía chân trời qua ô cửa sổ.

Nàng Sao chổi cầm một cái vợt bắt cá, hào hứng hỏi với giọng đầy phấn khích: ``Ayame-chan, bọn chị đang định đi săn khủng long đấy. Muốn tham gia không?''

\img{7-4h1}

Kuon nhìn cái vợt mà cô đang cầm, nghiêng đầu tỏ vẻ khó tin, giọng cậu đầy hoài nghi: ``Ý em là đi tìm xương khủng long đúng không? Với cả cái vợt đó thì làm được trò trống gì?''

Chưa kịp để hai người phản ứng, đằng sau họ, một bộ ba nghịch ngợm gồm nàng Matsuri, Haato và Fubuki đang nhảy nhót loạn cả lên, khiến cậu Tổng phải vội vàng đứng dậy can thiệp. Cậu lao về phía họ với vẻ mặt đầy bất lực: ``Này này, đừng có nhảy lên bàn! Mat-chan, đừng có xách váy Haachama lên! Còn Fubu-chan, đừng có giả tiếng khủng long nữa!''

Sau một hồi vật lộn với ba kẻ gây rối, cậu cuối cùng cũng tạm giữ chân được họ, thở dốc vì mệt. Nhưng chỉ trong tích tắc, họ đã lại phấn khích hét lên với giọng đầy nhiệt tình: ``\textbf{KHỦNG LONG!}''

Và thế là vòng lặp hỗn loạn lại tiếp tục, làm cậu Tổng vừa đuổi theo vừa cố gắng giữ họ trật tự, mồ hôi lấm tấm trên trán.

Nhưng rồi giữa mớ hỗn độn đó, Subaru bất ngờ chỉ tay về phía Ayame, người vẫn không hề rời mắt khỏi cuốn vở, giọng cô đầy vẻ bênh vực: ``Yên nào! Ayame-chan đang bận học để chuẩn bị cho kỳ thi lấy bằng Máy tính (cầm tay) điện tử đấy!''

\img{7-4h2}

``Toán bấm máy à? À, SAT chứ gì?'' cậu Tổng khẽ dừng lại, hỏi gọn với giọng đầy tò mò.

Nghe vậy, nàng Oni cuối cùng cũng đặt bút xuống, nhướng mày nhìn sang nàng Vịt với vẻ mặt không mấy hài lòng. Cô rướn người lên bàn, hạ giọng như thể sắp sửa sửa lưng đối phương: ``Ý của em là bài kiểm tra lấy bằng Hướng dẫn Trà đạo Nhật Bản cơ!''

Kuon nhíu mày, giọng cậu đầy bối rối: ``Bài kiểm tra gì nghe lạ thế?''

Nhưng chưa kịp hiểu chuyện gì, nàng Vịt lập tức bật lại với giọng đầy ngạc nhiên: ``Hả? Bài kiểm tra Hướng dẫn viên du lịch Internet á?''

``Bài kiểm tra Chuyên gia phục vụ và thử nếm khăn tắm!'' nàng Oni đáp lại, mặt bừng bừng tức giận vì đối phương không hiểu ý, hai tay cô nắm chặt thành đấm.

Kuon đơ mất mấy giây với những tên gọi lạ hoắc này, mắt cậu mở to với vẻ không thể tin nổi: ``\dots Cái củ quẹc gì vậy?'' Cậu sau đó thở dài, ngán ngẩm nói với giọng của một người đã bỏ cuộc: ``Thế này thì mấy đứa nên ôn bài kiểm tra `Phương trình Ngữ pháp và Cú pháp câu văn tiếng Nhật' đi, vì anh còn chả hiểu mấy đứa đang nói gì cả!''

Cậu vừa dứt câu, ba kẻ phá hoại đằng sau lại tiếp tục màn đồng thanh quen thuộc với giọng đầy nhiệt huyết: ``\textbf{Hai, ba! Khủng long!!}''

``Lại còn cả ba cái tên quỷ này nữa!?'' cậu Tổng quay ngoắt lại, mặt méo xệch, tay cậu chỉ về phía họ với vẻ bất lực. ``Làm ơn cư xử bình thường giúp anh cái!''

\img{7-4h3}

Suisei chỉ cười nhàn nhã với sự hỗn loạn này, giọng cô đầy vẻ triết lý: ``À\dots\ Đây đích thực là \hll\ mà chúng ta hay biết.'' Kuon đáp lại trong lúc vẫn đang ``úp sọt'' ba cô gái xuống bãi cát, giọng cậu đầy bất đắc dĩ: ``Dù anh ghét phải nói điều này, nhưng anh phải đồng ý\dots''

Cậu ôm trán, thở dài, thầm nghĩ với vẻ mặt đầy cam chịu: ``\textit{Mình đã làm gì để bị vây quanh bởi một đám `khủng long' dở hơi thế này chứ!?}''

\newpage

\subsection*{\phantomsection\label{ep-S13.IX}}
\addcontentsline{toc}{subsection}{{\sen{-Histoire 9-}} {\caviar \textit{hologra} cần gì để bùng nổ?}}
\stpt{-Histoire 9-}\stcap{\textit{hologra} cần gì để bùng nổ?}
\begin{center}
{\caviar\fontsize{12pt}{12pt}\selectfont Nhân vật: \emp{kuon1} \emp{achan} \emp{sora}}
\end{center}

Trong căn phòng họp yên tĩnh, Kuon đang ngồi cùng với A-chan và Sora. Trên bàn, một loạt giấy tờ về kế hoạch sản xuất \textit{holo no Graffiti} (\textit{hologra}) được bày ra, nhưng không ai có vẻ hào hứng lắm, những khuôn mặt đều trầm ngâm.

``Có ai có ý kiến gì không?'' nàng Đeo kính hỏi với giọng đầy mong đợi.

Kuon giơ tay, và cô gật đầu gọi cậu đứng dậy phát biểu. ``Chị đã có kế hoạch gì cho \textit{hologra} sắp tới chưa?''

Đồng nghiệp của cậu chống cằm, nhìn xuống tập giấy trước mặt với vẻ mặt mệt mỏi, rồi thở dài: ``Coi bộ cũng nhiều lắm, nhưng phần lớn toàn là xàm xí không à\dots\ Chẳng có gì đặc sắc cả.''

Cậu nhướng mày, giọng cậu đầy vẻ bất ngờ: ``Thì đặc sản của nó chả phải là sự xàm xí sao?''

A-chan khẽ cười, gật gù với vẻ đồng tình: ``Biết là thế nhưng\dots\ nó vẫn thiếu cái gì ấy.''

Trong lúc hai người nhân viên đang vặn óc suy nghĩ, Sora lúc này giơ tay lên, rồi lên tiếng với giọng đầy nhẹ nhàng: ``Sự nhiệt tình của mọi người?'' Nàng Đeo kính lắc đầu với vẻ mặt trầm tư: ``Không, chị không nghĩ thế. Nhiệt tình thì đúng là có thừa, nhưng\dots''

Nàng Thần tượng nghiêng đầu suy nghĩ, đôi mắt cô mơ màng với vẻ đoán già đoán non, rồi thử đoán tiếp: ``Vậy thì là sự nguyên bản?'' Kuon chống tay lên bàn, nhún vai với vẻ mặt không mấy tin tưởng: ``Anh không nghĩ vậy đâu. Nếu không phải là Yamazu viết kịch bản thì cũng là các thành viên tự biên tự diễn mà. Hay là thiếu người?''

Nàng Đeo kính xua tay với vẻ mặt của một người đang nghĩ đến tương lai: ``Từ sau 2nd Fes thì chúng ta sẽ tuyển thêm thành viên mới thôi. Chắc không phải cái đó đâu.''

Cả ba người thở dài chán nản, không khí phòng họp trở nên trầm lặng, ai cũng đang cố gắng tìm ra điều còn thiếu sót, những ánh mắt chăm chú vào các tập tài liệu.

Kuon lầm bầm trong miệng, giọng như một lời thì thầm: ``\hl{Chả nhẽ nó lại thiếu sự nổi bật để cho người nước ngoài cảm thấy thích thú?}''

Đột nhiên đôi mắt A-chan sáng bừng lên như vừa được khai sáng, cô đập tay xuống bàn: ``\textbf{Đúng rồi! Chính nó!}''

``Dạ?'' cậu Tổng ngạc nhiên, chưa kịp hiểu chuyện gì đang xảy ra.

Nàng Đeo kính đập mạnh tay xuống bàn, đầy phấn khích, giọng cô như một nhà phát minh vừa tìm ra công thức vĩ đại: ``Từ trước đến nay ta chỉ dùng những trò đùa hay những thứ có sẵn trong nước thôi, vậy tại sao ta không lấy những trò vui hay những tác phẩm từ nước ngoài nhỉ?''

Sora cũng vỗ tay đồng tình với cô ấy, mắt sáng lên với vẻ hào hứng: ``Cũng có thể lắm! Như thế thì sẽ càng có nhiều người nữa biết về chúng ta!''

Nàng Đeo kính nắm tay lại đầy quyết tâm, giọng cô như một lời tuyên bố: ``Quá tuyệt! Cảm ơn cậu nha, Hikawa-san!''

Hai người hào hứng bàn bạc về những nội dung quốc tế có thể đưa vào \textit{hologra}, bỏ lại cậu Tổng ngồi đó với vẻ mặt đầy sốc, mắt cậu mở to vì ngạc nhiên.

Nhìn hai cô nàng vui vẻ bàn tán với nhau, cậu vừa tỏ vẻ ngạc nhiên vừa hãnh diện khi cậu có thể giúp ích cho cả công ty này theo một cách nào đó, một nụ cười khẽ hiện trên môi: ``\textit{Mình nói vớ nói vẩn thế cũng được nữa\dots\ Hay thật!}''

\newpage

\subsection*{\phantomsection\label{ep-S13.X}}
\addcontentsline{toc}{subsection}{{\sen{-Histoire 10-}} {\caviar Nhìn lại một năm xàm xí}}
\stpt{-Histoire 10-}\stcap{Nhìn lại một năm xàm xí}
\begin{center}
{\caviar\fontsize{12pt}{12pt}\selectfont Nhân vật: \emp{kuon1} \emp{ryuuhou1}}
\end{center}

\begin{center}
  {\textit{*Các nhân vật nói bằng tiếng Etsunan.}}
\end{center}

Kuon đang ngồi trong phòng chiếu của trường, cắm tai nghe và dán mắt vào màn hình máy tính. Những thước phim cũ của \textit{hologra} lần lượt hiện lên, khiến cậu bật cười không ngừng, tay cậu gõ nhẹ lên bàn theo từng cảnh hài.

Lúc này cửa phòng mở ra. Hai người bạn cùng lớp là Shunyo và Tokue bước vào, thấy cậu cười tít mắt nên cậu bạn ``Sadboy'' liền hỏi với giọng đầy tò mò: ``Đại tướng xem cái gì mà cười cành cạch thế?''

Cậu tháo tai nghe ra, quay sang đáp với giọng đầy vẻ hài hước: ``À, mấy cái thước phim lộn xào mà tôi quay cùng với nhóm thần tượng holo ấy mà.''

Cậu bạn Hiến máu nhướng mày ngạc nhiên, khoanh tay lại với vẻ mặt có phần nghi hoặc: ``Ý ông là cái nhóm như cái rạp xiếc ấy á?'' Kuon gật đầu chắc nịch, giọng cậu như một lời xác nhận không thể chối cãi: ``Không họ thì ai.''

Cậu bạn ``Sadboy'' tò mò hỏi tiếp, giọng đầy thích thú: ``Hỏi vui tí chứ trong số này ông thích tập nào nhất?'' Kuon ngồi nghĩ một lúc lâu, mắt cậu nhìn xa xăm, rồi nở một nụ cười gian xảo, đáp gọn lỏn với giọng đầy ẩn ý: ``Tập 'mì trôi ống nước'.''

Cậu bạn Hiến máu nhăn mặt, giọng đầy vẻ khó hiểu: ``Là cái bìu gì thế?'' Cậu chàng hay cười khoanh tay, nheo mắt với vẻ tò mò: ``Sao lại thích nó thế, Đại tướng?''

Kuon bật cười, vỗ bàn với giọng đầy hào hứng: ``Đơn giản. Đồ ăn miễn phí. Đồng thời còn bắt nạt được cả Haachama nữa.''

``Thế nó như nào?'' cả hai cậu bạn cùng lớp đồng thanh, giọng đầy tò mò. Kuon không trả lời mà chỉ hướng màn hình cho hai người cùng xem, giọng cậu như một người dẫn chương trình: ``Cứ xem đi thì biết.''

\ct

Ban đầu cả hai chỉ định xem qua một chút cho vui, nhưng chẳng biết thế nào mà bị cuốn theo cả loạt video. Những tập \textit{hologra} từ năm 2020 lần lượt trôi qua, hết tập này đến tập khác, tiếng cười vang lên không ngớt.

Thỉnh thoảng, Tokue lại cười sặc sụa, còn cậu chàng hay cười thì hết lắc đầu lại bĩu môi khi họ nhìn lại chính mình (dù ít khi xuất hiện) trong đó, những biểu cảm đầy sự thích thú.

``Hài thật đó!'' cậu bạn Hiến máu vừa ôm bụng cười vừa nói, tay cậu đập nhẹ vào đùi.

``Đúng chứ? Mấy cái câu chuyện xàm xí này làm tôi dứt hết cả óc ra để giải quyết đấy,'' Kuon đáp lại, giọng đầy vẻ tự hào.

Shunyo chậc lưỡi, lắc đầu với vẻ mặt có phần khó chịu: ``Nhưng sao tôi thấy nó nhảm dị vậy\dots?''

Cậu Tổng chỉ nhún vai, đáp gọn lỏn với giọng của một người đã quá quen: ``Thì \hll\ là vậy mà.''

Cả ba cùng phá lên cười. Nhưng khi tiếng cười vừa lắng xuống, họ bỗng nhận ra xung quanh mình bỗng nhiên vắng tanh, không còn bóng dáng ai trong căn phòng vốn yên tĩnh.

Cậu bạn Hiến máu đảo mắt nhìn quanh phòng, giọng đầy ngạc nhiên: ``Ơ\dots? Mọi người đi đâu rồi?''

Kuon quay ra nhìn cửa sổ. Bầu trời bên ngoài đã dần chuyển sang sắc tím nhạt của hoàng hôn, những ánh đèn đường bắt đầu lóe sáng, báo hiệu một ngày đã qua. Cậu giật mình, đứng bật dậy với vẻ mặt hốt hoảng: ``Vãi!? Đến giờ này rồi sao!?''

Cậu bạn ``Sadboy'' cũng hốt hoảng, vội vàng thu dọn đồ đạc: ``Thôi, anh em mình về nhanh kẻo trễ giờ xe!''

Không ai bảo ai, cả ba đứa con trai ba chân bốn cẳng thu dọn đồ đạc rồi cùng nhau chạy ào ra khỏi phòng. Nhưng vừa ra đến hành lang, họ đụng phải một cô bé đang đứng tựa vào tường, tay cầm điện thoại, vẻ mặt đầy bất mãn.

Ryuuhou đứng đó với chiếc áo khoác dày, nhìn anh trai với ánh mắt nửa trách móc nửa buồn cười: ``Onii-san, em chờ từ chiều đến giờ luôn đấy. Nói là đi học chiều rồi về, thế mà giờ mới ra. Em tưởng anh bỏ quên em ở nhà hay sao ấy.''

Kuon sững lại, nhìn em gái với vẻ ngạc nhiên: ``Ryuuhou? Sao em lên đây?''

Ryuuhou lắc đầu, giọng đầy bất lực: ``Thì mẹ bảo em lên đón anh về vì anh quên điện thoại ở nhà nên mẹ không liên lạc được. Em gọi cũng không bắt máy, đành phải lên tận trường tìm. Ngồi ngoài cổng từ bốn giờ chiều mà không thấy anh đâu, cuối cùng phải hỏi bác bảo vệ mới biết anh vẫn còn trong phòng chiếu.''

Shunyo nhìn Ryuuhou rồi quay sang Kuon, giọng đầy trêu chọc: ``Ồ, ra là có em gái tới đón. Thế mà bảo vội về nhà đấy nhé.''

Kuon lúng túng, đưa tay xoa đầu: ``Xin lỗi em, anh mải xem phim với bọn bạn quá.''

Ryuuhou thở dài, nhưng nụ cười đã nở trên môi: ``Thôi, về đi. Mà này, anh xem gì mà mê đến thế vậy?''

``À, mấy tập \textit{hologra} cũ ấy mà.''

Mắt cô em gái sáng lên, cô bé nhảy chồm tới: ``Ồ! Cho em xem với! Em cũng muốn xem mấy cảnh xàm xí của các chị ấy!''

Kuon nhìn em gái, rồi quay sang hai người bạn, mỉm cười: ``Thôi, về nhà rồi tính. Mà hai đứa có muốn đi xe chung không? Chắc bọn mình trễ xe buýt rồi.''

Shunyo gật đầu: ``Có vẻ thế rồi. Nhờ cậu vậy.''

Ryuuhou lập tức giơ tay: ``Được, em gọi taxi cho mọi người nhé. Nhưng mà anh phải hứa là về nhà sẽ cho em xem mấy tập đó đấy.''

Ông anh trai lắc đầu cười: ``Rồi rồi, chiều em mà.''

Cả ba đứa con trai và cô em gái cùng nhau rời khỏi trường, tiếng cười giòn tan vang lên trên con đường dưới ánh đèn vàng của hoàng hôn. Họ chuẩn bị trở về quê đón một cái Tết Dương 2021, năm Tân Sửu, khép lại một năm 2020 Canh Tý dương lịch đầy hỗn loạn nhưng cũng tràn ngập tiếng cười và những kỷ niệm không thể quên.

\end{document}